\section{模型检验、评价与改进}

\subsection{数值精度与解析解验证}

本文的计算过程主要包含三类近似：一是将二阶动力学方程转化为一阶状态方程后
进行数值积分，二是在稳态区间内对瞬时功率作时间平均，三是在阻尼参数空间内
进行最优值搜索。对于线性阻尼情形，方程具有线性常系数结构，可以利用解析解
对数值算法进行校验；对于非线性阻尼和垂荡--纵摇耦合情形，系统含速度幂函数、
三角函数和惯性耦合项，难以直接获得闭式解析解，因此需要依赖数值积分和优化。

本文采用如下精度口径：对运动响应，关注位移和速度时间序列的数值误差；对优化
结果，关注稳态平均功率和最优阻尼参数的稳定性。由于非线性模型难以给出严格的
全局解析误差上界，本文以线性可解析问题的逐时刻误差、线性功率解析值与数值值
的偏差，以及候选最优点复算结果共同确定数值精度。

以问题一的线性阻尼为例，令
\[
  \boldsymbol z=(z_b,z_z)^\mathsf{T},
\]
系统可写为
\begin{equation}
  M\ddot{\boldsymbol z}+C\dot{\boldsymbol z}+K\boldsymbol z
  =\boldsymbol f\cos\omega t .
  \label{eq:evaluation_linear_matrix}
\end{equation}
其稳态解可设为
\[
  \boldsymbol z_s(t)=
  \mathrm{Re}\left(\boldsymbol Z e^{i\omega t}\right),
\]
其中复振幅为
\begin{equation}
  \boldsymbol Z=
  \left(K-\omega^2M+i\omega C\right)^{-1}\boldsymbol f .
  \label{eq:evaluation_complex_amplitude}
\end{equation}
进一步将式~\eqref{eq:evaluation_linear_matrix} 化为一阶状态方程
\[
  \dot{\boldsymbol x}=A\boldsymbol x+\boldsymbol b\cos\omega t,
\]
可得到满足零初始条件的完整解析解
\begin{equation}
  \boldsymbol x(t)=
  \mathrm{Re}\left(\boldsymbol q e^{i\omega t}\right)
  +e^{At}\left[\boldsymbol x(0)-\mathrm{Re}(\boldsymbol q)\right],
  \quad
  \boldsymbol q=(i\omega I-A)^{-1}\boldsymbol b .
  \label{eq:evaluation_exact_state_solution}
\end{equation}
将式~\eqref{eq:evaluation_exact_state_solution} 与 Runge--Kutta 数值积分结果
逐时刻比较，在 $40$ 个波浪周期内，浮子位移最大绝对误差为
$1.07\times10^{-10}\,\mathrm{m}$，振子位移最大绝对误差为
$1.12\times10^{-10}\,\mathrm{m}$；在最后 $10$ 个周期内，两者最大绝对误差
分别降至 $1.68\times10^{-11}\,\mathrm{m}$ 和
$1.74\times10^{-11}\,\mathrm{m}$。相对于约 $0.4\,\mathrm{m}$ 的响应
幅值，最大相对误差约为 $10^{-10}$ 量级，说明数值积分误差远小于模型参数
和工程近似带来的误差。

对于问题二的线性阻尼优化，阻尼系数 $c$ 虽为决策变量，但每一个给定的
$c$ 仍对应一个线性受迫振动系统。因此可直接由复振幅得到稳态平均功率
\begin{equation}
  \bar P(c)=
  \frac{1}{2}c\omega^2|Z_b(c)-Z_z(c)|^2 .
  \label{eq:evaluation_linear_power}
\end{equation}
利用式~\eqref{eq:evaluation_linear_power} 作一维优化，得到
\[
  c^\ast=37193.81\,\mathrm{N\cdot s/m},
  \qquad
  \bar P_{\max}=229.33\,\mathrm{W}.
\]
数值积分优化得到的阻尼系数为
$c\approx37420.87\,\mathrm{N\cdot s/m}$，代入解析功率公式后得到
$\bar P=229.33\,\mathrm{W}$，与解析最优功率仅相差
$4.22\times10^{-3}\,\mathrm{W}$，相对差约为
$1.84\times10^{-5}$。因此，本文对位移响应保留到
$10^{-6}\,\mathrm{m}$、对平均功率保留到 $10^{-2}\,\mathrm{W}$ 是有
数值精度支撑的；阻尼系数则考虑到功率峰附近平台较平，报告到整数位或两位
小数均不会改变工程结论。

\subsection{模型适用范围检验}

模型评价不仅要关注数值解是否收敛，还要检验计算结果是否落在建模假设允许的
物理范围内。本文主要从稳态截取、浮子垂荡振幅和纵摇小角度三个方面进行检验。

首先，平均输出功率应在系统进入稳定受迫振动后计算。初始静止条件会引入暂态
响应，若将暂态区间纳入功率平均，容易使阻尼参数评价偏离长期运行状态。因此，
本文在优化和复算阶段均舍弃前期过渡过程，只对后段近周期响应计算平均功率。以
问题二线性阻尼情形为例，采用更长时间积分并取后段窗口平均时，平均功率收敛到
$229.33\,\mathrm{W}$，与复振幅解析稳态功率一致。这说明最终采用的功率评价
口径反映的是系统稳定受迫振动下的长期平均输出，而不是初始条件引起的暂态输出。

问题四中同样对稳态区间作了复核。该问题波浪周期为
\[
  T=3.17236\,\mathrm{s}.
\]
在最优阻尼参数下，计算时舍弃前 $25$ 个周期，即从 $79.31\,\mathrm{s}$ 后
开始统计稳态平均功率。比较 $25T$ 到 $35T$ 和 $35T$ 到 $45T$ 两个区间，
总平均功率分别为 $320.63\,\mathrm{W}$ 和 $319.61\,\mathrm{W}$，相对差约为
$0.32\%$。这表明在问题四的功率统计区间内，系统响应已基本进入稳定周期状态。
最终采用 $25T$ 到 $45T$ 的整体稳态区间计算，得到平均总功率
$319.22\,\mathrm{W}$。

其次，线性静水恢复力 $K_hz_b$ 隐含水线始终位于浮子圆柱段内，即水线面积
保持为 $\pi r_b^2$。由静水平衡可得浮子和振子的总排水体积为
\[
  V_0=\frac{m_b+m_z}{\rho}=7.12098\,\mathrm{m^3}.
\]
浮子下部圆锥体积为
\[
  V_{\mathrm{cone}}=\frac{1}{3}\pi r_b^2h_{\mathrm{cone}}
  =0.83776\,\mathrm{m^3}.
\]
因此平衡时圆柱段浸没高度为
\[
  h_0=\frac{V_0-V_{\mathrm{cone}}}{\pi r_b^2}
  =2.00001\,\mathrm{m}.
\]
题给圆柱高度为 $h_{\mathrm{cyl}}=3\,\mathrm{m}$，取 $z_b$ 向下为正，则
水线仍位于圆柱段内需要满足
\[
  -2.00001\le z_b\le 0.99999.
\]
若采用单一振幅阈值作保守检验，则
\[
  A_{\max}=0.99999\,\mathrm{m}.
\]
计算结果中，问题二浮子垂荡最大位移约为 $0.792\,\mathrm{m}$，问题三约为
$0.830\,\mathrm{m}$，问题四约为 $0.822\,\mathrm{m}$，均小于
$A_{\max}$。因此，题给参数和本文最优阻尼附近的响应没有使圆锥段露出水面，
也没有使水线进入圆柱顶部以上区域，线性静水恢复力的适用条件未被破坏。

最后，检验纵摇角是否属于小量。本文在第三、四问中使用了题给的线性纵摇附加
转动惯量、线性兴波阻尼矩和线性静水恢复力矩，其中浮子纵摇方程含有
$-K_\theta\theta_b-C_\theta\dot\theta_b$。这些水动力系数和恢复力矩均是在静平衡
姿态附近的小扰动意义下给出的，因此需要验证求得的纵摇响应没有进入大角度范围。
同时，本文在几何耦合项中仍保留 $\sin\phi$ 和 $\cos\phi$ 的非线性形式，并未
把三角函数整体替换为线性近似；小角度检验的作用，是说明线性纵摇水动力和恢复
力矩的使用是自洽的。

问题三中，整个计算区间内浮子和振子的最大纵摇角为
\[
  \theta_{\max}=0.05205\,\mathrm{rad}=2.98^\circ .
\]
若从后段稳态区间统计，最大纵摇角进一步降为
$0.04022\,\mathrm{rad}=2.30^\circ$。若在恢复力矩线性化或局部误差估计中采用
小角度近似，则由泰勒展开余项估计，
\[
  \frac{|\sin\theta-\theta|}{|\theta|}\lesssim\frac{\theta^2}{6},
  \qquad
  |1-\cos\theta|\lesssim\frac{\theta^2}{2}.
\]
取全区间最大角度 $\theta_{\max}=0.05205$，可得正弦近似相对误差不超过
$4.52\times10^{-4}$，余弦近似绝对误差不超过 $1.35\times10^{-3}$。
问题四的纵摇角更小，全区间最大纵摇角为
$0.04435\,\mathrm{rad}=2.54^\circ$；在 $25T$ 后的稳态统计区间内，最大
纵摇角为 $0.03160\,\mathrm{rad}=1.81^\circ$，此时正弦近似相对误差上界约为
$1.66\times10^{-4}$，余弦近似绝对误差上界约为 $4.99\times10^{-4}$。
因此，在题给激励和最优阻尼附近，纵摇响应始终保持小角度范围。

\subsection{优化结果稳定性分析}

阻尼参数优化的目标函数由受迫振动响应和功率积分共同决定。对于这类问题，最优
参数未必表现为尖锐的孤立点；若功率曲面在峰值附近较平，则一组参数都可能具有
近似相同的工程效果。因此，除了给出最优点，还需要分析最优结果的稳定性。

问题二线性阻尼情形可以由解析稳态功率函数直接判断最优区间。根据
式~\eqref{eq:evaluation_linear_power}，当
\[
  c\in[35561.29,38901.28]\,\mathrm{N\cdot s/m}
\]
时，稳态平均功率均不低于解析最优值的 $99.9\%$。这说明
$c=37193.81\,\mathrm{N\cdot s/m}$ 与数值优化得到的
$c=37420.87\,\mathrm{N\cdot s/m}$ 虽有差异，但对应平均功率几乎相同。
从工程角度看，线性阻尼参数可在上述区间内选择，而不必过度强调单点最优值。

问题二非线性阻尼优化得到
\[
  c=8.7065\times10^4,\qquad \alpha=0.3593,
\]
对应平均功率约为 $227.03\,\mathrm{W}$。与线性阻尼最优功率相比，非线性阻尼
在题给参数下带来的功率提升并不显著，说明该问题的主要优化空间仍来自阻尼强度
与系统固有响应之间的匹配，而非单纯依赖幂指数形式。由于非线性功率曲面存在
较平缓区域，幂指数和阻尼系数的数值不宜解释为绝对唯一的工程参数，而应理解为
在给定模型和频率下的一组有效匹配参数。

问题四中，最优阻尼参数为
\[
  c_1=59846.67\,\mathrm{N\cdot s/m},\qquad
  c_2=67472.34\,\mathrm{N\cdot m\cdot s}.
\]
稳态平均总功率为 $319.22\,\mathrm{W}$，其中垂荡阻尼贡献
$319.18\,\mathrm{W}$，旋转阻尼贡献仅 $0.047\,\mathrm{W}$。旋转阻尼功率
占总功率的比例约为
\[
  \frac{0.047}{319.22}=1.47\times10^{-4}.
\]
这与前述纵摇角较小的运动特征一致，也说明总功率主要由垂荡阻尼决定。因而在
问题四中，旋转阻尼最优值附近可能存在较大的等效选择空间；评价该结果时应更
重视总功率平台和垂荡阻尼匹配，而非将旋转阻尼的单点数值解释为唯一精确解。

\subsection{模型优点}

本文模型首先从受力和力矩平衡出发，变量定义和物理含义较为清晰。问题一、二
将系统抽象为浮子和振子的二自由度垂荡受迫振动模型，能够直接描述 PTO 弹簧和
阻尼器通过相对运动输出功率的机制；问题三、四进一步引入浮子纵摇、振子相对
转动以及振子沿浮子中轴滑移，使模型能够覆盖题目要求的主要运动自由度。

其次，本文没有完全依赖黑箱数值计算。在线性阻尼情形下，利用复振幅法和矩阵
指数形式给出解析解，并用解析结果对数值积分和线性功率优化进行验证；在非线性
阻尼和垂荡--纵摇耦合情形下，再采用数值积分方法求解。这样的处理既保留了线性
模型的精确性，又保证了后续复杂模型的可计算性。

再次，优化方法兼顾了全局搜索和局部精修。对阻尼参数先作粗扫描以观察功率曲面
形状和峰值区域，再用局部优化方法精修候选最优点，并通过高精度复算确认平均
功率。这种策略降低了局部优化对初值的依赖，也便于发现功率曲面较平坦时存在的
近似最优参数区间。

\subsection{模型局限与改进方向}

本文模型仍有若干局限。第一，模型采用固定静水面和线性静水恢复力，虽然前文已
验证本文计算结果没有超出水线位于圆柱段的适用范围，但当波浪激励进一步增大时，
水面相对浮子的位置和湿表面形状会发生明显变化，此时应引入瞬时水面高度和非线性
浮力模型。

第二，本文仅针对题给规则波和固定频率参数进行优化。实际海况往往由随机波谱或
多频波组成，最优阻尼可能随频率分布、波高和海况统计特征改变。后续可引入真实
波谱，对不同海况下的平均捕能效率进行评价，并寻找更具鲁棒性的阻尼参数。

第三，本文将 PTO 阻尼耗散功率等同于可输出功率，未考虑机械传动、电机转换、
控制策略和能量存储环节的损失。在工程应用中，应进一步建立 PTO 系统效率模型，
区分机械吸收功率和最终电功率。

第四，本文主要考虑平面内垂荡和纵摇运动，忽略横荡、横摇和艏摇等三维自由度。
对于实际海上装置，波浪入射方向和结构姿态可能导致更多自由度耦合。后续可将
本文的受力和力矩平衡框架推广到六自由度模型，以提高复杂海况下的适用性。
